\documentclass[11pt]{article}

\usepackage[margin=1in]{geometry}
\usepackage{amsmath,amssymb}
\usepackage{booktabs,tabularx}
\usepackage{enumitem}
\usepackage[dvipsnames]{xcolor}
\usepackage[most]{tcolorbox}
\usepackage{microtype}
\usepackage[hidelinks]{hyperref}

\definecolor{accent}{HTML}{244A73}
\definecolor{paleblue}{HTML}{F1F6FA}
\newcolumntype{Y}{>{\raggedright\arraybackslash}X}
\newtcolorbox{keybox}[1][Key idea]{
  colback=paleblue,colframe=accent,fonttitle=\bfseries,
  title=#1,boxrule=0.6pt,arc=1.5mm,left=2mm,right=2mm,top=1mm,bottom=1mm}
\newcommand{\Pp}{\operatorname{P}}
\newcommand{\E}{\operatorname{E}}
\newcommand{\Var}{\operatorname{Var}}
\newcommand{\SD}{\operatorname{SD}}
\newcommand{\Bin}{\operatorname{Bin}}
\newcommand{\N}{\operatorname{N}}
\newcommand{\Unif}{\operatorname{U}}
\newcommand{\dd}{\,\mathrm{d}}

\hypersetup{pdftitle={Ten Essential Concepts in Probability},
  pdfauthor={Study Notes}}
\setlist{nosep}
\setcounter{tocdepth}{2}

\title{\textbf{Ten Essential Concepts in Probability}\\[0.4em]
  \large A rigorous study guide for competition preparation and review}
\author{}
\date{}

\begin{document}
\maketitle

\begin{abstract}
Probability measures how likely an event is to occur.  For any event $A$, its
probability satisfies $0\leq \Pp(A)\leq 1$, where $0$ means impossible and $1$
means certain.  Probability questions begin by identifying the random experiment,
the sample space $S$ of elementary outcomes, and the event $A\subseteq S$ of
interest.  These notes develop ten essential ideas: estimating probability from
experiments, computing probabilities from finite models, using sets and
conditional probability, counting with permutations and combinations, connecting
events to random variables, and recognizing continuous, binomial, and geometric
probability models.  The emphasis is on choosing a valid model, checking its
assumptions, calculating exactly when possible, and interpreting the answer.
\end{abstract}

\tableofcontents
\newpage

\section{Experimental Probability}\label{sec:experimental}
A \emph{random experiment} is a process whose outcome is uncertain before it is
performed.  One performance of the experiment is a \emph{trial}.  An
\emph{outcome} is a possible result; an \emph{elementary outcome} is one
indivisible result in the chosen model.  The \emph{sample space} $S$ is the set
of elementary outcomes under consideration, and an \emph{event} is a subset of
$S$.

\begin{keybox}[Vocabulary example]
For one roll of a fair die, the experiment is rolling the die once, one trial is
one roll, the elementary outcomes are $1,2,3,4,5,6$,
\[
S=\{1,2,3,4,5,6\},\qquad A=\{1,3,5\}
\]
is the event ``an odd result.''
\end{keybox}

\emph{Experimental probability} is a probability estimated from observed data.
If event $A$ occurs in $n(A)$ of $N$ trials, then the experimental probability,
or relative frequency, is
\[
\widehat{\Pp}(A)=\frac{n(A)}{N},\qquad N>0.
\]
This formula is valid after the trials have actually been performed.

\subsection{Worked example: flipping a coin}
Suppose a coin is flipped $10$ times and tails occurs $3$ times.  Then
\[
\widehat{\Pp}(\text{tails})=\frac{3}{10}=0.3=30\%.
\]
This means that tails occurred in $30\%$ of the flips in this particular
experiment.  It does \emph{not} mean that the coin's theoretical probability of
tails is $0.3$.

\subsection{Experimental versus theoretical probability}
For a fair coin, the theoretical probability of tails is $1/2=0.5=50\%$.
Experimental probability is a sample proportion, so repeated experiments with
the same number of trials can produce different relative frequencies.  For
independent repeated trials with the same event probability $p=\Pp(A)$, the
relative frequency converges toward $p$ in the long run.  Individual experiments
can still fluctuate, and adding more trials does not guarantee monotonic
improvement after every trial; larger sample sizes generally reduce typical
fluctuation.

\begin{table}[ht]
\centering
\caption{Experimental and theoretical probability.}\label{tab:exp-theory}
\begin{tabularx}{\textwidth}{@{}lYY@{}}
\toprule
Type & Based on & Main caution \\
\midrule
Experimental probability & Observed outcomes from trials & May vary from one experiment to another \\
Theoretical probability & A mathematical model of the sample space & Requires correct assumptions about the model \\
\bottomrule
\end{tabularx}
\end{table}

Simulation uses artificial repetitions of a model to estimate a probability that
is difficult to calculate directly.  For example, to estimate the probability of
getting at least one six in five die rolls, one could repeatedly simulate five
fair rolls, record whether at least one six appeared, and compute the relative
frequency.  Simulation estimates the model probability; it does not replace the
need to check whether the model is appropriate.

\begin{keybox}
Experimental probability answers, ``What happened in these trials?''  Theoretical
probability answers, ``What should happen according to the model?''
\end{keybox}

\section{Theoretical Probability}\label{sec:theoretical}
\emph{Theoretical probability} uses a mathematical description of the experiment
to compute the probability of an event.  In a finite model, probabilities obey
\[
\Pp(A)\geq0,\qquad \Pp(S)=1,
\]
and if $A\cap B=\varnothing$, then
\[
\Pp(A\cup B)=\Pp(A)+\Pp(B).
\]
Consequences include
\[
\Pp(\varnothing)=0,\qquad \Pp(A^c)=1-\Pp(A),\qquad
A\subseteq B\Longrightarrow \Pp(A)\leq \Pp(B).
\]
These are finite-level probability rules; no advanced set theory is needed here.

For a finite sample space $S$ whose elementary outcomes are equally likely, the
probability of event $A\subseteq S$ is
\[
\Pp(A)=\frac{|A|}{|S|}.
\]
Here $|A|$ is the number of favorable outcomes and $|S|$ is the total number of
possible outcomes.  This favorable-outcomes formula is valid only when all
elementary outcomes in $S$ are equally likely.

\subsection{Worked examples}
In these examples, assume fair dice, a well-shuffled standard deck, and random
selection.
\begin{enumerate}[label=\textbf{\arabic*.}]
\item Rolling a four on a fair six-sided die:
\[
\Pp(\text{four})=\frac{1}{6}.
\]
\item Rolling an odd number:
\[
\Pp(\text{odd})=\frac{3}{6}=\frac{1}{2}=0.5=50\%.
\]
\item Drawing the seven of hearts from a standard $52$-card deck:
\[
\Pp(\text{seven of hearts})=\frac{1}{52}.
\]
\item Drawing a king:
\[
\Pp(\text{king})=\frac{4}{52}=\frac{1}{13}.
\]
\item Drawing a red marble from a bag with $3$ red, $2$ blue, and $5$ green
marbles.  The elementary outcomes are the individual marbles, not merely the
three color names, so
\[
\Pp(\text{red})=\frac{3}{3+2+5}=\frac{3}{10}=0.3.
\]
\end{enumerate}

\subsection{Complements and weighted finite models}
The \emph{complement} $A^c$ is the event that $A$ does not occur.  For the
marble bag above,
\[
\Pp(\text{not blue})=1-\Pp(\text{blue})
=1-\frac{2}{10}=\frac{8}{10}=\frac{4}{5}=0.8=80\%.
\]
The notation $A'$ is also sometimes used for a complement, but $A^c$ is more
standard in many probability texts.

If elementary outcomes are not equally likely, add their probabilities instead
of counting them.  For a finite sample space,
\[
\Pp(A)=\sum_{\omega\in A}\Pp(\{\omega\}),
\]
where each elementary probability is nonnegative and all elementary probabilities
sum to $1$.

\begin{table}[ht]
\centering
\caption{Constructed weighted finite sample space.}\label{tab:weighted}
\begin{tabular}{@{}ccc@{}}
\toprule
Outcome & $a$ & $b$ & $c$ \\
\midrule
Probability & $0.2$ & $0.5$ & $0.3$ \\
\bottomrule
\end{tabular}
\end{table}

In this constructed model, the event $\{a,c\}$ has probability
$0.2+0.3=0.5$, not $2/3$.  Counting outcomes alone fails because the elementary
outcomes are not equally likely.

\begin{keybox}[Common error]
Do not use $\Pp(A)=|A|/|S|$ unless the elementary outcomes are equally likely.
A biased coin, a nonuniform spinner, or a weighted table requires probability
weights.
\end{keybox}

\section{Probability Using Sets}\label{sec:sets}
\emph{Probability using sets} describes events as subsets of a sample space and
uses set operations to combine them.  The main notation is as follows.
\begin{itemize}
\item $S$: the sample space, or set of all possible outcomes.
\item Event: any subset of $S$.
\item $\varnothing$: the empty event, which contains no outcomes.
\item $A^c$: the complement of $A$, meaning outcomes not in $A$.
\item $A\cap B$: the intersection, meaning outcomes in both $A$ and $B$.
\item $A\cup B$: the union, meaning outcomes in $A$ or $B$ or both.
\item Mutually exclusive events: events with no overlap, so $A\cap B=\varnothing$.
\end{itemize}

The word \emph{and} usually signals an intersection.  The word \emph{or} in
probability is inclusive: $A$ or $B$ means $A$, $B$, or both.

\begin{table}[ht]
\centering
\caption{Language translated into event notation.}\label{tab:set-language}
\begin{tabularx}{\textwidth}{@{}lY@{}}
\toprule
Phrase & Event notation \\
\midrule
$A$ and $B$ & $A\cap B$ \\
$A$ or $B$; at least one of $A,B$ & $A\cup B$ \\
Only $A$ & $A\cap B^c$ \\
Neither $A$ nor $B$; none of $A,B$ & $(A\cup B)^c=A^c\cap B^c$ \\
Exactly one of $A,B$ & $(A\cap B^c)\cup(A^c\cap B)$ \\
\bottomrule
\end{tabularx}
\end{table}

De Morgan's laws describe complements of combined events:
\[
(A\cup B)^c=A^c\cap B^c,
\qquad
(A\cap B)^c=A^c\cup B^c.
\]
A \emph{partition} of $S$ is a collection of mutually exclusive events whose
union is all of $S$.  Partitions are useful because they let us split a
probability into disjoint cases.

\subsection{The addition rule}
For any two events $A$ and $B$,
\[
\Pp(A\cup B)=\Pp(A)+\Pp(B)-\Pp(A\cap B).
\]
The subtraction removes the overlap, which was counted once in $\Pp(A)$ and
again in $\Pp(B)$.  If $A$ and $B$ are mutually exclusive, then
$\Pp(A\cap B)=0$, so $\Pp(A\cup B)=\Pp(A)+\Pp(B)$.

\begin{table}[ht]
\centering
\caption{Venn-style regions for two events.}\label{tab:venn-regions}
\begin{tabularx}{\textwidth}{@{}lY@{}}
\toprule
Region & Description \\
\midrule
$A\cap B$ & Outcomes in both events \\
$A\cap B^c$ & Outcomes in $A$ only \\
$A^c\cap B$ & Outcomes in $B$ only \\
$(A\cup B)^c$ & Outcomes in neither event \\
\bottomrule
\end{tabularx}
\end{table}

\subsection{Worked example: a sports survey}
A survey has $100$ respondents.  The verified information is: $8$ respondents
like both hockey and basketball; $48$ like basketball; $42$ like soccer; and
$15$ like both basketball and soccer.  Let $B$ be ``likes basketball'' and $C$
be ``likes soccer.''  The probability that a randomly selected respondent likes
basketball or soccer is
\begin{align*}
\Pp(B\cup C)&=\frac{48+42-15}{100}\\
&=\frac{75}{100}=\frac{3}{4}=0.75.
\end{align*}
Thus $75\%$ of respondents are known to be in the basketball-or-soccer union.
The value $8/100=0.08$ is the probability of liking hockey and basketball.  No
unreported regions of the original Venn diagram are needed for these conclusions.

\begin{keybox}
When using a Venn diagram, add each region exactly once.  For a union, subtract
the intersection if the totals for two sets have already counted it twice.
\end{keybox}

\section{Conditional Probability}\label{sec:conditional}
\emph{Conditional probability} is the probability of one event under the
condition that another event has occurred.  The probability of $B$ given $A$ is
\[
\Pp(B\mid A)=\frac{\Pp(A\cap B)}{\Pp(A)},\qquad \Pp(A)>0.
\]
For a finite equally counted sample, this can be written as
\[
\Pp(B\mid A)=\frac{|A\cap B|}{|A|},\qquad |A|>0.
\]
Conditioning restricts the relevant sample space from all of $S$ to the event
$A$.

\subsection{Worked example: school preference}
The following complete table is a constructed classroom example.
\begin{table}[ht]
\centering
\caption{Constructed school-preference table.}\label{tab:school}
\begin{tabular}{@{}lccc@{}}
\toprule
 & Likes school $L$ & Does not like school $L^c$ & Total \\
\midrule
Female $F$ & $10$ & $4$ & $14$ \\
Not female $F^c$ & $8$ & $8$ & $16$ \\
\midrule
Total & $18$ & $12$ & $30$ \\
\bottomrule
\end{tabular}
\end{table}

Then
\[
\Pp(L\mid F)=\frac{|L\cap F|}{|F|}=\frac{10}{14}=\frac{5}{7}\approx0.7143,
\]
while
\[
\Pp(F\mid L)=\frac{|F\cap L|}{|L|}=\frac{10}{18}=\frac{5}{9}\approx0.5556.
\]
The reversed conditionals are different because their denominators are
different.  These are associations within the data, not causal claims.

\subsection{Total probability and Bayes' theorem}
If $B_1,\ldots,B_m$ form a partition of $S$ and each relevant denominator is
positive, then
\[
\Pp(A)=\sum_{i=1}^{m}\Pp(B_i)\Pp(A\mid B_i).
\]
Bayes' theorem reverses a conditional probability:
\[
\Pp(B_j\mid A)=
\frac{\Pp(B_j)\Pp(A\mid B_j)}
{\sum_{i=1}^{m}\Pp(B_i)\Pp(A\mid B_i)},\qquad \Pp(A)>0.
\]

\paragraph{Base-rate example.}
A screening test is applied in a population where $2\%$ have a condition.  The
test is positive for $90\%$ of people with the condition and for $5\%$ of people
without it.  Let $D$ be ``has the condition'' and $+$ be ``positive test.''
Then
\begin{align*}
\Pp(+)&=\Pp(D)\Pp(+\mid D)+\Pp(D^c)\Pp(+\mid D^c)\\
&=(0.02)(0.90)+(0.98)(0.05)=0.067.
\end{align*}
Thus
\[
\Pp(D\mid +)=\frac{(0.02)(0.90)}{0.067}
=\frac{18}{67}\approx0.2687.
\]
The base rate $\Pp(D)=0.02$ matters; $\Pp(+\mid D)$ and $\Pp(D\mid +)$ answer
different questions.

\begin{keybox}[Common error]
In $\Pp(B\mid A)$, the denominator is the size or probability of the condition
$A$, not the entire original sample space.
\end{keybox}

\section{The Multiplication Law}\label{sec:multiplication}
The \emph{multiplication law} computes the probability that events occur
together.  Whenever the conditional probabilities are defined,
\[
\Pp(A\cap B)=\Pp(A)\Pp(B\mid A)=\Pp(B)\Pp(A\mid B).
\]
Events $A$ and $B$ are \emph{independent} if and only if
\[
\Pp(A\cap B)=\Pp(A)\Pp(B).
\]
When $\Pp(A)>0$ and $\Pp(B)>0$, this is equivalent to
\[
\Pp(B\mid A)=\Pp(B),\qquad \Pp(A\mid B)=\Pp(A).
\]

For sequential events,
\[
\Pp(A_1\cap\cdots\cap A_m)=\Pp(A_1)\Pp(A_2\mid A_1)\cdots
\Pp(A_m\mid A_1\cap\cdots\cap A_{m-1}).
\]

\begin{table}[ht]
\centering
\caption{Independent, dependent, and mutually exclusive events.}\label{tab:independent-dependent}
\begin{tabularx}{\textwidth}{@{}lYY@{}}
\toprule
Relationship & Meaning & Important distinction \\
\midrule
Independent & Product rule $\Pp(A\cap B)=\Pp(A)\Pp(B)$ holds & Can overlap \\
Dependent & The product rule with unconditional probabilities does not hold & Often occurs in sampling without replacement \\
Mutually exclusive & Cannot occur together & If both have positive probability, they are not independent \\
\bottomrule
\end{tabularx}
\end{table}

\subsection{Worked examples}
\paragraph{A die and a coin.} Rolling a fair die and flipping a fair coin are
independent.  Therefore,
\[
\Pp(\text{three and tails})=\frac{1}{6}\cdot\frac{1}{2}=\frac{1}{12}.
\]

\paragraph{Two kings without replacement.} Drawing two kings consecutively from
a standard deck without replacement is dependent.  The first draw has $4$ kings
among $52$ cards.  If the first card is a king, then $3$ kings remain among
$51$ cards, so
\[
\Pp(\text{two kings})=\frac{4}{52}\cdot\frac{3}{51}
=\frac{12}{2652}=\frac{1}{221}.
\]
The numerator and denominator change because the first card is not put back.

\subsection{Replacement and branch tables}
A branch table is a text version of a probability tree: multiply along a path
and add the paths that satisfy the event.  Suppose a bag contains $4$ red and
$6$ blue marbles.

\begin{table}[ht]
\centering
\caption{Two draws with replacement.}\label{tab:with-replacement}
\begin{tabular}{@{}cccc@{}}
\toprule
First draw & Second draw & Path probability & Interpretation \\
\midrule
$R$ & $R$ & $\frac{4}{10}\cdot\frac{4}{10}=\frac{4}{25}$ & probabilities unchanged \\
$R$ & $B$ & $\frac{4}{10}\cdot\frac{6}{10}=\frac{6}{25}$ & probabilities unchanged \\
$B$ & $R$ & $\frac{6}{10}\cdot\frac{4}{10}=\frac{6}{25}$ & probabilities unchanged \\
$B$ & $B$ & $\frac{6}{10}\cdot\frac{6}{10}=\frac{9}{25}$ & probabilities unchanged \\
\bottomrule
\end{tabular}
\end{table}

With replacement, $\Pp(\text{one red and one blue})=6/25+6/25=12/25$.

\begin{table}[ht]
\centering
\caption{Two draws without replacement.}\label{tab:without-replacement}
\begin{tabular}{@{}cccc@{}}
\toprule
First draw & Second draw & Path probability & Interpretation \\
\midrule
$R$ & $R$ & $\frac{4}{10}\cdot\frac{3}{9}=\frac{2}{15}$ & probabilities change \\
$R$ & $B$ & $\frac{4}{10}\cdot\frac{6}{9}=\frac{4}{15}$ & probabilities change \\
$B$ & $R$ & $\frac{6}{10}\cdot\frac{4}{9}=\frac{4}{15}$ & probabilities change \\
$B$ & $B$ & $\frac{6}{10}\cdot\frac{5}{9}=\frac{1}{3}$ & probabilities change \\
\bottomrule
\end{tabular}
\end{table}

Without replacement, $\Pp(\text{one red and one blue})=4/15+4/15=8/15$.

\begin{keybox}
Use multiplication for ``and'' along a path, but decide first whether later
probabilities must be conditional.  Add disjoint successful paths.
\end{keybox}

\section{Permutations}\label{sec:permutations}
The \emph{fundamental counting principle} says that if a process has $m_1$
choices at the first stage, $m_2$ choices at the second, and so on, then the
total number of possible outcomes is
\[
m_1m_2\cdots m_r.
\]
For example, a lunch with $3$ sandwich choices, $2$ drink choices, and $4$ side
choices can be selected in $3\cdot2\cdot4=24$ ways.

A \emph{permutation} is an ordered arrangement of objects.  For a nonnegative
integer $n$, factorial notation means
\[
n!=n(n-1)(n-2)\cdots2\cdot1,\qquad 0!=1.
\]
For $0\leq r\leq n$, the number of ordered arrangements of $r$ distinct objects
chosen from $n$ distinct objects is
\[
{}_nP_r=P(n,r)=\frac{n!}{(n-r)!}.
\]
This formula applies when order matters and repetition is not allowed.

If order matters and repetition is allowed, then each of the $r$ positions has
$n$ choices, so there are
\[
n^r
\]
ordered selections.

\subsection{Worked examples}
The six arrangements of $A,B,C$ are
\[
ABC,\ ACB,\ BAC,\ BCA,\ CAB,\ CBA.
\]
For first-, second-, and third-place finishes among $10$ runners,
\[
{}_{10}P_3=\frac{10!}{7!}=10\cdot9\cdot8=720.
\]

For a lock code of length $3$ using numbers $0$ through $59$:
\begin{itemize}
\item if repetition is allowed, there are $60^3=216{,}000$ possible codes;
\item if the three numbers must be distinct, there are
\[
{}_{60}P_3=60\cdot59\cdot58=205{,}320
\]
possible codes.
\end{itemize}
If only one distinct-number code is correct, then
\[
\Pp(\text{correct first guess})=\frac{1}{205{,}320}.
\]
Check: $60\cdot59\cdot58=205{,}320$.

\begin{keybox}
A permutation question asks, ``How many ordered lists can be made?''  If swapping
two selected objects creates a different outcome, order matters.
\end{keybox}

\section{Combinations}\label{sec:combinations}
A \emph{combination} is a selection in which order does not matter.  For
$0\leq r\leq n$, the number of ways to choose $r$ objects from $n$ distinct
objects is
\[
\binom{n}{r}=\frac{n!}{r!(n-r)!}.
\]
The extra factor $r!$ in the denominator removes the different orders of the
same selected group.  Equivalently,
\[
{}_nP_r=r!\binom{n}{r},
\]
because one may first choose the $r$ objects and then arrange the selected
objects in $r!$ orders.

\subsection{Worked examples}
The number of groups of three from five people is
\[
\binom{5}{3}=\frac{5!}{3!2!}=10.
\]
Now suppose a team of $6$ is chosen from $7$ children and $11$ adults.  The
sample space contains
\[
\binom{18}{6}=18{,}564
\]
teams.  To form a favorable team with exactly one child, choose $1$ of the $7$
children and $5$ of the $11$ adults:
\[
\binom{7}{1}\binom{11}{5}=7\cdot462=3{,}234.
\]
Therefore
\[
\Pp(\text{exactly one child})=
\frac{\binom{7}{1}\binom{11}{5}}{\binom{18}{6}}
=\frac{3234}{18564}=\frac{77}{442}\approx0.1742.
\]

\subsection{Complementary counting}
Sometimes it is easier to count what does \emph{not} happen:
\[
\Pp(\text{at least one})=1-\Pp(\text{none}).
\]
For example, the probability of at least one six in two fair die rolls is
\[
1-\Pp(\text{no sixes})=1-\left(\frac56\right)^2
=1-\frac{25}{36}=\frac{11}{36}\approx0.3056.
\]

\begin{table}[ht]
\centering
\caption{Permutations and combinations.}\label{tab:perm-comb}
\begin{tabularx}{\textwidth}{@{}lYY@{}}
\toprule
Counting method & Use when & Formula \\
\midrule
Ordered with repetition & Order matters, repetition allowed & $n^r$ \\
Permutation & Order matters, no repetition & ${}_nP_r=\dfrac{n!}{(n-r)!}$ \\
Combination & Order does not matter & $\binom{n}{r}=\dfrac{n!}{r!(n-r)!}$ \\
\bottomrule
\end{tabularx}
\end{table}

\begin{keybox}[Connection to binomial probabilities]
In a binomial probability, $\binom{n}{k}$ counts the possible locations of the
$k$ successes among the $n$ trials.
\end{keybox}

\section{Continuous Probability Distributions}\label{sec:continuous}
Before named distributions, connect events to numerical values.  A \emph{random
variable} $X$ assigns a number to each outcome of an experiment.  Its
\emph{support} is the set of values it can take.  A discrete random variable has
possible values that can be listed or counted; a continuous random variable has
possible values occupying intervals.

For three fair coin flips, let $X$ be the number of heads.  Then the support is
$\{0,1,2,3\}$; for example, $X(\text{HTH})=2$.  A \emph{probability mass
function} (PMF) is
\[
p_X(k)=\Pp(X=k),\qquad p_X(k)\geq0,
\qquad \sum_k p_X(k)=1.
\]
For two fair coin flips, the PMF of $X=$ number of heads is:
\begin{table}[ht]
\centering
\caption{PMF for heads in two fair flips.}\label{tab:pmf-heads}
\begin{tabular}{@{}cccc@{}}
\toprule
$k$ & $0$ & $1$ & $2$ \\
\midrule
$\Pp(X=k)$ & $\frac14$ & $\frac12$ & $\frac14$ \\
\bottomrule
\end{tabular}
\end{table}
The probabilities are nonnegative and sum to $1/4+1/2+1/4=1$.

The \emph{cumulative distribution function} (CDF) of any random variable is
\[
F(x)=\Pp(X\leq x).
\]
Thus $\Pp(X>x)=1-F(x)$.  For integer-valued discrete variables, endpoints matter:
for example, $\Pp(X\geq k)=1-F(k-1)$.

For a discrete random variable,
\[
\E(X)=\sum_x x\Pp(X=x)
\]
is its expected value, interpreted as a long-run average.  It need not be a
possible individual observation.  Variance and standard deviation measure spread:
\[
\Var(X)=\E\!\left[(X-\E(X))^2\right],
\qquad
\SD(X)=\sqrt{\Var(X)}.
\]

A \emph{continuous random variable} can take any value in an interval of real
numbers.  A \emph{probability density function} (PDF) $f(x)$ assigns density,
not point probability, and satisfies
\[
f(x)\geq0,
\qquad
\int_{-\infty}^{\infty} f(x)\dd x=1.
\]
The integral represents area under the density curve; no integration techniques
are needed in these notes.  For a continuous random variable $X$ with density
$f$,
\[
\Pp(a\leq X\leq b)=\int_a^b f(x)\dd x.
\]
Thus probabilities are areas under the curve, not curve heights.  Also,
$\Pp(X=x)=0$ for an individual value, so endpoints do not affect interval
probabilities: $\Pp(a\leq X\leq b)=\Pp(a<X<b)$.

\subsection{Uniform distribution}
A uniform distribution on an interval gives equal density to equal lengths.  We
write
\[
X\sim \Unif(a,b),\qquad f(x)=\frac{1}{b-a}\quad\text{for }a\leq x\leq b.
\]
If $X\sim\Unif(0,10)$, then
\[
\Pp(2<X<7)=\frac{7-2}{10-0}=\frac12=0.5.
\]
The shaded region in a graph would be the interval from $2$ to $7$ under a flat
rectangle of height $1/10$.

\subsection{The normal distribution}
A normal model is appropriate only when context or assumptions support an
approximately symmetric, bell-shaped distribution; not every continuous variable
or every bell-shaped curve is normal.  We write
$X\sim \N(\mu,\sigma^2)$ when $X$ is normal with mean $\mu$, variance
$\sigma^2$, and standard deviation $\sigma>0$.  Approximately $68\%$, $95\%$,
and $99.7\%$ of the area lie within $1$, $2$, and $3$ standard deviations of
the mean, respectively, under a normal model.

The standard normal distribution is
\[
Z\sim \N(0,1),
\]
with mean $0$, variance $1$, and standard deviation $1$.  Its CDF is
\[
\Phi(z)=\Pp(Z\leq z).
\]
To standardize a normal value, use
\[
Z=\frac{X-\mu}{\sigma}\sim\N(0,1).
\]
Useful identities are
\[
\Pp(Z>z)=1-\Phi(z),\qquad
\Pp(a<Z<b)=\Phi(b)-\Phi(a),\qquad
\Phi(-z)=1-\Phi(z).
\]

\subsection{Worked examples: normal-model calculations}
Suppose
\[
X\sim \N(68,10^2),\qquad \text{and find }\Pp(X>76).
\]
The standardized value is
\[
z=\frac{76-68}{10}=0.8.
\]
Using a standard normal table or a normal CDF,
\[
\Phi(0.8)\approx0.7881.
\]
Therefore
\[
\Pp(X>76)=1-\Phi(0.8)\approx1-0.7881=0.2119.
\]
Under the stated normal-model assumption, the probability is approximately
$0.2119$, or $21.19\%$.

For a between-values calculation, suppose $X\sim\N(100,15^2)$.  Then
\[
\Pp(85<X<115)=\Pp\left(-1<Z<1\right)
=\Phi(1)-\Phi(-1)\approx0.8413-0.1587=0.6826.
\]
For percentiles, if $\Phi(z_q)=q$, then
\[
x_q=\mu+\sigma z_q.
\]
For example, the $90$th percentile of $\N(100,15^2)$ is approximately
$100+15(1.2816)=119.224$, because $\Phi(1.2816)\approx0.9000$.

\begin{table}[ht]
\centering
\caption{Discrete and continuous distributions.}\label{tab:discrete-continuous}
\begin{tabularx}{\textwidth}{@{}lYY@{}}
\toprule
Type & Probability tool & Endpoint behavior \\
\midrule
Discrete & Probability mass function & Individual values can have positive probability \\
Continuous & Probability density function and area & Individual values have probability $0$ \\
\bottomrule
\end{tabularx}
\end{table}

\begin{keybox}[Common error]
A density height is not the probability of an exact value, and a $z$-score is not
a probability.  Probabilities come from CDF values or areas over intervals.
\end{keybox}

\section{Binomial Probability Distributions}\label{sec:binomial}
A \emph{binomial probability distribution} models the number of successes in a
fixed number of Bernoulli trials.  The conditions are:
\begin{enumerate}
\item a fixed number $n$ of trials;
\item two categorized outcomes per trial, called success and failure;
\item independent trials;
\item a constant success probability $p$ on every trial.
\end{enumerate}
If $X$ is the number of successes, then
\[
X\sim\Bin(n,p)
\]
and the PMF is
\[
\Pp(X=k)=\binom{n}{k}p^k(1-p)^{n-k},\qquad k=0,1,\ldots,n.
\]
A particular sequence with $k$ successes and $n-k$ failures has probability
$p^k(1-p)^{n-k}$.  The factor $\binom{n}{k}$ counts the possible locations of
the $k$ successes among the $n$ trials.  The expected value, variance, and
standard deviation are
\[
\E(X)=np,
\qquad
\Var(X)=np(1-p),
\qquad
\SD(X)=\sqrt{np(1-p)}.
\]

For cumulative probabilities,
\[
\Pp(X\leq r)=\sum_{k=0}^{r}\Pp(X=k),
\]
and complements often simplify calculations, such as
\[
\Pp(X\geq1)=1-\Pp(X=0).
\]

\begin{keybox}[Binomial checklist]
Before using a binomial model, ask: fixed $n$? two categorized outcomes?
independent trials? constant $p$?  Sampling without replacement generally
changes $p$ and creates dependence, so it is not exactly binomial; a large
population may justify an approximation, but that is a modeling judgment.
\end{keybox}

\subsection{Worked examples}
\paragraph{Rolling threes.} Roll a fair six-sided die four times and define
success as rolling a three.  Then
\[
X\sim\Bin\left(4,\frac16\right).
\]
The complete probability distribution is:
\begin{table}[ht]
\centering
\caption{Number of threes in four rolls.}\label{tab:binomial}
\begin{tabular}{@{}ccc@{}}
\toprule
$k$ & Exact probability $\Pp(X=k)$ & Decimal approximation \\
\midrule
$0$ & $\left(\frac56\right)^4=\frac{625}{1296}$ & $0.4823$ \\
$1$ & $\binom41\left(\frac16\right)\left(\frac56\right)^3=\frac{500}{1296}$ & $0.3858$ \\
$2$ & $\binom42\left(\frac16\right)^2\left(\frac56\right)^2=\frac{150}{1296}$ & $0.1157$ \\
$3$ & $\binom43\left(\frac16\right)^3\left(\frac56\right)=\frac{20}{1296}$ & $0.0154$ \\
$4$ & $\left(\frac16\right)^4=\frac{1}{1296}$ & $0.0008$ \\
\bottomrule
\end{tabular}
\end{table}
The exact probabilities sum to
\[
\frac{625+500+150+20+1}{1296}=\frac{1296}{1296}=1.
\]
The expected number of threes is
\[
\E(X)=4\left(\frac16\right)=\frac23\approx0.667.
\]
The variance is
\[
\Var(X)=4\left(\frac16\right)\left(\frac56\right)=\frac{5}{9},
\qquad \SD(X)=\frac{\sqrt5}{3}\approx0.7454.
\]
The expected value is a long-run average over many repetitions of the four-roll
experiment; it need not be a possible observed count in one experiment.

\paragraph{At least one six.} If a fair die is rolled $5$ times and $X$ is the
number of sixes, then $X\sim\Bin(5,1/6)$ and
\[
\Pp(X\geq1)=1-\Pp(X=0)=1-\left(\frac56\right)^5
=1-\frac{3125}{7776}=\frac{4651}{7776}\approx0.5981.
\]

\paragraph{A failed binomial checklist.} Drawing $5$ cards from a deck without
replacement and counting aces is not exactly binomial: the number of trials is
fixed and there are two categorized outcomes, but the draws are dependent and
the success probability changes after each draw.

\section{Geometric Probability Distributions}\label{sec:geometric}
A \emph{geometric probability distribution} models the trial number on which the
first success occurs.  The conditions are:
\begin{enumerate}
\item independent repeated Bernoulli trials;
\item two categorized outcomes per trial;
\item constant success probability $p$;
\item trials continue until the first success.
\end{enumerate}
If $X$ is the trial number of the first success, then
\[
\Pp(X=k)=(1-p)^{k-1}p,
\qquad k=1,2,\ldots.
\]
The expected waiting time and variance are
\[
\E(X)=\frac{1}{p},
\qquad
\Var(X)=\frac{1-p}{p^2}.
\]
Tail and cumulative probabilities are often easiest by complements:
\[
\Pp(X>k)=(1-p)^k,
\qquad
\Pp(X\leq k)=1-(1-p)^k.
\]
Here $\Pp(X>k)$ means no success occurs in the first $k$ trials.

\subsection{Worked example: rolling doubles}
Roll two fair six-sided dice repeatedly until doubles appear.  There are $36$
equally likely ordered outcomes and $6$ doubles, so
\[
p=\Pp(\text{doubles})=\frac{6}{36}=\frac16.
\]
The probability that the first doubles occur on trial $4$ is
\[
\Pp(X=4)=\left(\frac56\right)^3\left(\frac16\right)
=\frac{125}{1296}\approx0.09645.
\]
For fewer than three trials,
\begin{align*}
\Pp(X<3)&=\Pp(X=1)+\Pp(X=2)\\
&=\frac16+\frac56\cdot\frac16
=\frac{6}{36}+\frac{5}{36}
=\frac{11}{36}\approx0.3056.
\end{align*}
The probability that doubles do not occur in the first five trials is
\[
\Pp(X>5)=\left(\frac56\right)^5=\frac{3125}{7776}\approx0.4019.
\]
The expected waiting time is
\[
\E(X)=\frac{1}{1/6}=6,
\]
and the variance is
\[
\Var(X)=\frac{1-1/6}{(1/6)^2}=30.
\]
On average, it takes six rolls of the pair of dice to get doubles.

\subsection{Binomial or geometric?}
For fair coin flips, ``exactly two heads in five flips'' is binomial:
\[
\binom52\left(\frac12\right)^2\left(\frac12\right)^3=\frac{10}{32}=\frac{5}{16}.
\]
The question fixes five trials and counts successes.  By contrast, ``the first
head occurs on flip five'' is geometric:
\[
\left(\frac12\right)^4\left(\frac12\right)=\frac{1}{32}.
\]
The process continues until the first success.

Situations that are not geometric include drawing from a finite deck without
replacement, any process with changing success probability, and an experiment
that stops after a fixed number of trials regardless of whether success occurs.

\begin{keybox}[Memoryless property]
For a geometric random variable, the chance of waiting more than $t$ additional
trials does not depend on already having waited $s$ trials:
\[
\Pp(X>s+t\mid X>s)=\Pp(X>t).
\]
This property relies on independent trials with constant $p$.
\end{keybox}

\begin{table}[ht]
\centering
\caption{Binomial and geometric models.}\label{tab:binom-geometric}
\begin{tabularx}{\textwidth}{@{}lYY@{}}
\toprule
Model & Random variable & Experiment stops when \\
\midrule
Binomial & Number of successes in $n$ trials & The fixed $n$ trials are completed \\
Geometric & Trial number of first success & The first success occurs \\
\bottomrule
\end{tabularx}
\end{table}

\section{Final Synthesis}\label{sec:synthesis}
The progression of ideas is
\[
\text{experiment}\longrightarrow
\text{sample space and events}\longrightarrow
\text{counting and conditioning}\longrightarrow
\text{random variable}\longrightarrow
\text{distribution}.
\]
Experimental probability estimates likelihood from data, while theoretical
probability uses a model.  Set notation clarifies intersections, inclusive
unions, complements, and partitions.  Conditional probability changes the sample
space, and the multiplication law turns sequential ``and'' questions into
products, using conditional probabilities when events are dependent.

Permutations and combinations support probability by counting sample spaces and
event spaces.  Use permutations when order matters and combinations when only
the selected group matters.  Random variables turn outcomes into numbers.
Continuous models assign probability through areas under density curves.  The
binomial and geometric distributions are discrete PMFs for repeated Bernoulli
trials, but they answer different questions: ``How many successes in a fixed
number of trials?'' versus ``How long until the first success?''

\begin{table}[ht]
\centering
\caption{Model-selection checklist.}\label{tab:model-selection}
\begin{tabularx}{\textwidth}{@{}lY@{}}
\toprule
Question to ask & Likely tool \\
\midrule
Are elementary outcomes equally likely? & Use favorable outcomes over total outcomes. \\
Are elementary outcomes weighted unequally? & Add the weights of favorable outcomes. \\
Is there a condition or restricted sample space? & Use conditional probability. \\
Are events sequentially dependent? & Use conditional multiplication or a branch table. \\
Does order matter? Is repetition allowed? & Choose among $n^r$, permutations, and combinations. \\
Is the random variable discrete or continuous? & Use a PMF/CDF or a PDF/area model. \\
Is there a fixed number of Bernoulli trials? & Consider a binomial model after checking assumptions. \\
Does the process continue until the first success? & Consider a geometric model after checking assumptions. \\
Are the assumptions actually satisfied? & Revise or reject the model if assumptions fail. \\
\bottomrule
\end{tabularx}
\end{table}

\section{Formula Reference}\label{sec:formulas}
\begin{table}[ht]
\centering
\caption{Core formulas from the ten topics.}\label{tab:formulas}
\small
\begin{tabularx}{\textwidth}{@{}lY@{}}
\toprule
Topic & Formula and conditions \\
\midrule
Foundations & $\Pp(A)\ge0$, $\Pp(S)=1$, $\Pp(\varnothing)=0$; if $A\cap B=\varnothing$, then $\Pp(A\cup B)=\Pp(A)+\Pp(B)$ \\
Consequences & $\Pp(A^c)=1-\Pp(A)$; if $A\subseteq B$, then $\Pp(A)\le\Pp(B)$ \\
Experimental probability & $\widehat{\Pp}(A)=n(A)/N$, after $N>0$ observed trials \\
Equally likely outcomes & $\Pp(A)=|A|/|S|$, finite equally likely sample space \\
Weighted finite space & $\Pp(A)=\sum_{\omega\in A}\Pp(\{\omega\})$, elementary probabilities nonnegative and summing to $1$ \\
Union & $\Pp(A\cup B)=\Pp(A)+\Pp(B)-\Pp(A\cap B)$ \\
De Morgan's laws & $(A\cup B)^c=A^c\cap B^c$ and $(A\cap B)^c=A^c\cup B^c$ \\
Conditional probability & $\Pp(B\mid A)=\Pp(A\cap B)/\Pp(A)$, $\Pp(A)>0$ \\
Multiplication law & $\Pp(A\cap B)=\Pp(A)\Pp(B\mid A)=\Pp(B)\Pp(A\mid B)$ when defined \\
Chain rule & $\Pp(A_1\cap\cdots\cap A_m)=\Pp(A_1)\Pp(A_2\mid A_1)\cdots\Pp(A_m\mid A_1\cap\cdots\cap A_{m-1})$ \\
Independence & $\Pp(A\cap B)=\Pp(A)\Pp(B)$; if denominators are positive, $\Pp(B\mid A)=\Pp(B)$ and $\Pp(A\mid B)=\Pp(A)$ \\
Total probability & If $B_1,\ldots,B_m$ partition $S$, then $\Pp(A)=\sum_i\Pp(B_i)\Pp(A\mid B_i)$ \\
Bayes' theorem & $\Pp(B_j\mid A)=\dfrac{\Pp(B_j)\Pp(A\mid B_j)}{\sum_i\Pp(B_i)\Pp(A\mid B_i)}$, $\Pp(A)>0$ \\
Counting principle & $m_1m_2\cdots m_r$ choices over $r$ stages \\
Ordered with repetition & $n^r$ ordered selections when each position has $n$ choices \\
Permutations & ${}_nP_r=n!/(n-r)!$, order matters, no repetition \\
Combinations & $\binom{n}{r}=n!/[r!(n-r)!]$, order does not matter; ${}_nP_r=r!\binom nr$ \\
PMF and CDF & $p_X(x)=\Pp(X=x)$, $p_X(x)\ge0$, $\sum_xp_X(x)=1$; $F(x)=\Pp(X\le x)$; $P(X>x)=1-F(x)$ \\
Expected value and spread & $\E(X)=\sum_xx\Pp(X=x)$ for discrete $X$; $\Var(X)=\E[(X-\E(X))^2]$; $\SD(X)=\sqrt{\Var(X)}$ \\
Continuous density & $f(x)\ge0$, $\int_{-\infty}^{\infty}f(x)\dd x=1$, $\Pp(a\le X\le b)=\int_a^bf(x)\dd x$, and $\Pp(X=x)=0$ \\
Uniform & $X\sim\Unif(a,b)$ has $f(x)=1/(b-a)$ for $a\le x\le b$ \\
Normal & $\Phi(z)=\Pp(Z\le z)$; if $X\sim\N(\mu,\sigma^2)$ with $\sigma>0$, then $Z=(X-\mu)/\sigma\sim\N(0,1)$ \\
Normal CDF identities & $\Pp(Z>z)=1-\Phi(z)$, $\Pp(a<Z<b)=\Phi(b)-\Phi(a)$, $\Phi(-z)=1-\Phi(z)$ \\
Binomial & $\Pp(X=k)=\binom nkp^k(1-p)^{n-k}$, $k=0,\ldots,n$, for fixed independent Bernoulli trials with constant $p$ \\
Binomial mean/spread & $\E(X)=np$, $\Var(X)=np(1-p)$, $\SD(X)=\sqrt{np(1-p)}$; $\Pp(X\le r)=\sum_{k=0}^r\Pp(X=k)$; $\Pp(X\ge1)=1-\Pp(X=0)$ \\
Geometric & $\Pp(X=k)=(1-p)^{k-1}p$, $k=1,2,\ldots$, for the trial number of the first success \\
Geometric tail/mean/spread & $\Pp(X>k)=(1-p)^k$, $\Pp(X\le k)=1-(1-p)^k$, $\E(X)=1/p$, $\Var(X)=(1-p)/p^2$ \\
\bottomrule
\end{tabularx}
\end{table}

\section{Mixed Review Exercises}\label{sec:review}
\begin{enumerate}
\item A coin is flipped $20$ times and heads appears $11$ times.  Find the
experimental probability of heads.
\item A card is drawn from a standard $52$-card deck.  Find the probability that
it is not a king.
\item In a class of $30$ students, $18$ play soccer, $12$ play basketball, and
$7$ play both.  Let $C$ be ``plays soccer'' and $B$ be ``plays basketball.''
What is $\Pp(C\cup B)$?
\item A bag contains $4$ red and $6$ blue marbles.  Two marbles are drawn without
replacement.  Find the probability that both are red.
\item How many ways can gold, silver, and bronze medals be awarded among $12$
competitors if no competitor can receive more than one medal?
\item How many committees of $4$ can be formed from $9$ people?
\item Suppose $X\sim\N(50,8^2)$, where $8^2$ is the variance.  Express
$\Pp(X<42)$ using the standard normal CDF $\Phi$.
\item A fair die is rolled $5$ times.  Let $X$ be the number of sixes.  Find
$\Pp(X=2)$.
\item A fair coin is flipped until the first head.  Find the probability that the
first head occurs on the third flip.
\item Explain whether the situation in Exercise 8 is binomial or geometric.
\item A weighted spinner has outcomes $a,b,c$ with probabilities $0.2,0.5,0.3$.
Find $\Pp(\{a,c\})$ and explain why $2/3$ is not correct.
\item In a constructed survey of $40$ students, $24$ study music, $18$ study art,
and $10$ study both.  Translate ``exactly one of music and art'' into event
notation and find its probability.
\item A table has $12$ athletes who like chess, $8$ athletes who do not like
chess, $6$ non-athletes who like chess, and $14$ non-athletes who do not like
chess.  Let $A$ be ``athlete'' and $H$ be ``likes chess.''  Find $\Pp(H\mid A)$
and $\Pp(A\mid H)$.
\item A population has $30\%$ users of a service.  A survey response is positive
for $80\%$ of users and $20\%$ of nonusers.  Find the probability of a positive
response and then $\Pp(\text{user}\mid\text{positive})$.
\item A bag contains $3$ red and $2$ blue marbles.  Compare $\Pp(RR)$ for two
draws with replacement and without replacement.
\item A code has three positions and each position can use one of the digits
$0$--$9$.  How many codes are possible if repetition is allowed?  How many if
repetition is prohibited?
\item For two fair coin flips, let $X$ be the number of heads.  Give the support
of $X$ and its PMF.
\item If $X\sim\Unif(0,10)$, find $\Pp(2<X<7)$.
\item For $X\sim\Bin(5,1/6)$, find $\Pp(X\ge1)$.
\item For a geometric random variable with $p=1/6$, find $\Pp(X>5)$.
\item Decide whether drawing $4$ cards without replacement and counting hearts is
exactly binomial.  Explain using the binomial checklist.
\item Choose the most appropriate method for each task and justify briefly:
\begin{enumerate}[label=(\alph*)]
\item probability of a password when order matters and repetition is allowed;
\item probability of a positive test after base rates and conditional test rates
are given;
\item probability that a continuous waiting time lies between two numbers;
\item probability that the first success occurs on trial $6$.
\end{enumerate}
\end{enumerate}

\subsection*{Solutions}
\begin{enumerate}
\item $\widehat{\Pp}(\text{heads})=11/20=0.55=55\%$.
\item $\Pp(\text{not king})=1-4/52=48/52=12/13$.
\item Use the union rule:
\[
\Pp(C\cup B)=\frac{18+12-7}{30}=\frac{23}{30}\approx0.7667.
\]
\item The draws are dependent:
\[
\Pp(\text{two red})=\frac{4}{10}\cdot\frac{3}{9}=\frac{2}{15}\approx0.1333.
\]
\item This is a permutation because medal order matters:
\[
{}_{12}P_3=12\cdot11\cdot10=1320.
\]
\item This is a combination because committee order does not matter:
\[
\binom{9}{4}=126.
\]
\item Standardize using standard deviation $\sigma=8$:
\[
z=\frac{42-50}{8}=-1,
\qquad \Pp(X<42)=\Phi(-1)\approx0.1587.
\]
\item Here $X\sim\Bin(5,1/6)$, so
\[
\Pp(X=2)=\binom52\left(\frac16\right)^2\left(\frac56\right)^3
=\frac{625}{3888}\approx0.1608.
\]
\item The first two flips must be tails and the third must be heads:
\[
\Pp(X=3)=\left(\frac12\right)^2\left(\frac12\right)=\frac18=0.125.
\]
\item Exercise 8 is binomial because there is a fixed number of trials, $5$, and
we count the number of successes.  Exercise 9 is geometric because trials
continue until the first success.
\item Add the weights of the favorable elementary outcomes:
\[
\Pp(\{a,c\})=0.2+0.3=0.5.
\]
The answer $2/3$ incorrectly treats $a,b,c$ as equally likely.
\item Exactly one is $(M\cap A^c)\cup(M^c\cap A)$, where $M$ is music and $A$ is
art.  The number studying exactly one is
\[
(24-10)+(18-10)=14+8=22,
\]
so the probability is $22/40=11/20=0.55$.
\item The total number of athletes is $12+8=20$, and the total number who like
chess is $12+6=18$.  Therefore
\[
\Pp(H\mid A)=\frac{12}{20}=\frac35=0.6,
\qquad
\Pp(A\mid H)=\frac{12}{18}=\frac23\approx0.6667.
\]
The reversed conditionals differ because their denominators differ.
\item Let $U$ be user and $+$ be positive.  By total probability,
\[
\Pp(+)=(0.30)(0.80)+(0.70)(0.20)=0.24+0.14=0.38.
\]
Bayes' theorem gives
\[
\Pp(U\mid +)=\frac{(0.30)(0.80)}{0.38}=\frac{0.24}{0.38}
=\frac{12}{19}\approx0.6316.
\]
\item With replacement,
\[
\Pp(RR)=\frac35\cdot\frac35=\frac{9}{25}=0.36.
\]
Without replacement,
\[
\Pp(RR)=\frac35\cdot\frac24=\frac{3}{10}=0.3.
\]
The second branch probability changes without replacement.
\item With repetition allowed, there are $10^3=1000$ codes.  Without repetition,
there are
\[
{}_{10}P_3=10\cdot9\cdot8=720
\]
codes.
\item The support is $\{0,1,2\}$, and
\[
\Pp(X=0)=\frac14,
\qquad
\Pp(X=1)=\frac12,
\qquad
\Pp(X=2)=\frac14.
\]
These probabilities sum to $1$.
\item For a uniform distribution on $[0,10]$,
\[
\Pp(2<X<7)=\frac{7-2}{10-0}=\frac12=0.5.
\]
\item Use the complement:
\[
\Pp(X\ge1)=1-\Pp(X=0)=1-\left(\frac56\right)^5
=\frac{4651}{7776}\approx0.5981.
\]
\item For the first-success convention,
\[
\Pp(X>5)=\left(1-\frac16\right)^5=\left(\frac56\right)^5
=\frac{3125}{7776}\approx0.4019.
\]
\item It is not exactly binomial.  The number of draws is fixed and each draw can
be classified as heart or not heart, but without replacement the draws are
dependent and the probability of a heart changes after each draw.
\item The appropriate methods are:
\begin{enumerate}[label=(\alph*)]
\item ordered counting with repetition, using $n^r$;
\item total probability and Bayes' theorem, because base rates and conditional
rates are given;
\item a continuous density or CDF calculation, because the variable is continuous
and the event is an interval;
\item a geometric model, because the question asks for the first success on a
specified trial.
\end{enumerate}
\end{enumerate}

\end{document}
